\documentclass[a4paper,12pt]{article}
\usepackage[utf8]{inputenc}
\usepackage{amsmath, amssymb, amsthm} % For math symbols and environments
\usepackage{graphicx} % For including images
\usepackage{hyperref} % For hyperlinks
\usepackage{fancyhdr} % For header and footer
\usepackage{geometry} % For adjusting page dimensions
\usepackage{cite} % For handling references
\geometry{margin=1in}

% Header and footer
\pagestyle{fancy}
\fancyhf{}
\fancyhead[L]{Sravanth Chowdary Potluri}
\fancyhead[C]{MLIA PS1}
\fancyhead[R]{\thepage}
\fancyfoot[L]{CS6501 - 005}
\fancyfoot[R]{\today}

% Title Page
\title{Machine Learning In Image Analysis PS1}
\author{Sravanth Chowdary Potluri \\ und5uv \\ CS6501 - 005}
\date{\today}

\begin{document}

\maketitle

% Question 1
\section{Question 1}
Prove the following properties of convolution. For all continuous functions $f$, $g$, and $h$, the following axioms hold. \textbf{Make sure to lay out all key steps of your proof for full credit}.

\begin{itemize}
    \item \textbf{Associativity:} $(f \ast g) \ast h = f \ast (g \ast h)$
    \item \textbf{Distributivity:} $f \ast (g + h) = f \ast g + f \ast h$
    \item \textbf{Differentiation Rule:} $(f \ast g)' = f' \ast g = f \ast g'$
    \item \textbf{Convolution Theorem:} $\mathcal{F}(g \ast h) = \mathcal{F}(g) \mathcal{F}(h)$, where $\mathcal{F}$ denotes the Fourier transform.
\end{itemize}



\section*{Solution}
Please note that the functions and notations used here are slightly different to as explained in class, However they are accurate and true.

Let \( f, g, \) and \( h \) be continuous functions. The convolution of two functions is defined as:
\[
(f \ast g)(t) = \int_{-\infty}^{\infty} f(\tau) g(t - \tau) \, d\tau.
\]
We will prove the following properties using the integral definition of convolution.

---

\subsection*{1. Associativity}
\textbf{Statement:} \( (f \ast g) \ast h = f \ast (g \ast h) \).

\textbf{Proof:}

First, compute \( (f \ast g) \ast h \):
\[
\begin{aligned}
((f \ast g) \ast h)(t) &= \int_{-\infty}^{\infty} (f \ast g)(\tau) h(t - \tau) \, d\tau \\
&= \int_{-\infty}^{\infty} \left[ \int_{-\infty}^{\infty} f(u) g(\tau - u) \, du \right] h(t - \tau) \, d\tau.
\end{aligned}
\]

Switch the order of integration (by Fubini's theorem):
\[
\begin{aligned}
((f \ast g) \ast h)(t) &= \int_{-\infty}^{\infty} f(u) \left[ \int_{-\infty}^{\infty} g(\tau - u) h(t - \tau) \, d\tau \right] du.
\end{aligned}
\]

Let \( s = \tau - u \), so \( \tau = s + u \) and \( d\tau = ds \). Then,
\[
\begin{aligned}
((f \ast g) \ast h)(t) &= \int_{-\infty}^{\infty} f(u) \left[ \int_{-\infty}^{\infty} g(s) h(t - s - u) \, ds \right] du \\
&= \int_{-\infty}^{\infty} f(u) \left[ (g \ast h)(t - u) \right] du \\
&= (f \ast (g \ast h))(t).
\end{aligned}
\]

\textbf{Conclusion:} \( (f \ast g) \ast h = f \ast (g \ast h) \).

---

\subsection*{2. Distributivity}
\textbf{Statement:} \( f \ast (g + h) = f \ast g + f \ast h \).

\textbf{Proof:}

Compute \( f \ast (g + h) \):
\[
\begin{aligned}
(f \ast (g + h))(t) &= \int_{-\infty}^{\infty} f(\tau) [g + h](t - \tau) \, d\tau \\
&= \int_{-\infty}^{\infty} f(\tau) [g(t - \tau) + h(t - \tau)] \, d\tau \\
&= \int_{-\infty}^{\infty} f(\tau) g(t - \tau) \, d\tau + \int_{-\infty}^{\infty} f(\tau) h(t - \tau) \, d\tau \\
&= (f \ast g)(t) + (f \ast h)(t).
\end{aligned}
\]

\textbf{Conclusion:} \( f \ast (g + h) = f \ast g + f \ast h \).

---

\subsection*{3. Differentiation Rule}
\textbf{Statement:} \( (f \ast g)' = f' \ast g = f \ast g' \).

\textbf{Proof:}

\textbf{First Part: \( (f \ast g)' = f \ast g' \)}

Compute the derivative of the convolution:
\[
\begin{aligned}
(f \ast g)'(t) &= \frac{d}{dt} \int_{-\infty}^{\infty} f(\tau) g(t - \tau) \, d\tau \\
&= \int_{-\infty}^{\infty} f(\tau) \frac{d}{dt} g(t - \tau) \, d\tau \\
&= \int_{-\infty}^{\infty} f(\tau) g'(t - \tau) \, d\tau \\
&= (f \ast g')(t).
\end{aligned}
\]

\textbf{Second Part: \( (f \ast g)' = f' \ast g \)}

Alternatively, consider:
\[
\begin{aligned}
(f \ast g)(t) &= \int_{-\infty}^{\infty} f(\tau) g(t - \tau) \, d\tau \\
&= \int_{-\infty}^{\infty} f(t - s) g(s) \, ds \quad (\text{Let } s = t - \tau) \\
&= \int_{-\infty}^{\infty} f(t - s) g(s) \, ds.
\end{aligned}
\]

Differentiate with respect to \( t \):
\[
\begin{aligned}
(f \ast g)'(t) &= \int_{-\infty}^{\infty} \frac{d}{dt} f(t - s) g(s) \, ds \\
&= \int_{-\infty}^{\infty} f'(t - s) g(s) \, ds \\
&= (f' \ast g)(t).
\end{aligned}
\]

\textbf{Conclusion:} \( (f \ast g)' = f' \ast g = f \ast g' \).

---

\subsection*{4. Convolution Theorem}
\textbf{Statement:} \( \mathcal{F}(g \ast h) = \mathcal{F}(g) \cdot \mathcal{F}(h) \), where \( \mathcal{F} \) denotes the Fourier transform.

\textbf{Proof:}

definition of the Fourier transform:
\[
\mathcal{F}(f)(\omega) = \int_{-\infty}^{\infty} f(t) e^{-i \omega t} \, dt.
\]

Compute \( \mathcal{F}(g \ast h)(\omega) \):
\[
\begin{aligned}
\mathcal{F}(g \ast h)(\omega) &= \int_{-\infty}^{\infty} (g \ast h)(t) e^{-i \omega t} \, dt \\
&= \int_{-\infty}^{\infty} \left[ \int_{-\infty}^{\infty} g(\tau) h(t - \tau) \, d\tau \right] e^{-i \omega t} \, dt \\
&= \int_{-\infty}^{\infty} g(\tau) \left[ \int_{-\infty}^{\infty} h(t - \tau) e^{-i \omega t} \, dt \right] d\tau.
\end{aligned}
\]

Let \( s = t - \tau \), so \( t = s + \tau \) and \( dt = ds \):
\[
\begin{aligned}
\mathcal{F}(g \ast h)(\omega) &= \int_{-\infty}^{\infty} g(\tau) \left[ \int_{-\infty}^{\infty} h(s) e^{-i \omega (s + \tau)} \, ds \right] d\tau \\
&= \int_{-\infty}^{\infty} g(\tau) e^{-i \omega \tau} \left[ \int_{-\infty}^{\infty} h(s) e^{-i \omega s} \, ds \right] d\tau \\
&= \left[ \int_{-\infty}^{\infty} g(\tau) e^{-i \omega \tau} \, d\tau \right] \left[ \int_{-\infty}^{\infty} h(s) e^{-i \omega s} \, ds \right] \\
&= \mathcal{F}(g)(\omega) \cdot \mathcal{F}(h)(\omega).
\end{aligned}
\]

\textbf{Conclusion:} \( \mathcal{F}(g \ast h) = \mathcal{F}(g) \cdot \mathcal{F}(h) \).

\section{Question 2}

Frequency Smoothing

\begin{enumerate}
    \item[(a)] Compute the Fourier transform of the given image \texttt{lenaNoise.PNG} by using \texttt{numpy.fft.fft2} in Python, and then center the low frequencies (e.g., by using \texttt{fftshift}).
    
    \item[(b)] Keep a different number of low frequencies (e.g., 72, 152, 312, and the full dimension), but set all other high frequencies to 0.
    
    \item[(c)] Reconstruct the original image (using \texttt{ifft2}) by using the new generated frequencies in step (b).
\end{enumerate}

\section*{Solution}
Solution In Jupyter Notebook which is Also Attached After The Report

\section{Question 3}

\subsection{Problem Description}

Image denoising is the process of removing noise from a given noisy image to recover a clean image. Let \( f \in \mathbb{R}^{m \times n} \) be the observed noisy image, where each pixel value contains additive noise. The goal is to find a clean image \( u \in \mathbb{R}^{m \times n} \) that is as close as possible to the noisy image \( f \), while ensuring that the resulting image \( u \) is smooth and free from noise.

One common approach to achieve this is by minimizing an energy function that consists of two terms:
\begin{itemize}
    \item A data fidelity term, which ensures that \( u \) stays close to the noisy image \( f \).
    \item A regularization term, which imposes smoothness on \( u \) while preserving important features like edges.
\end{itemize}

\subsection{Objective Function}

The objective function to minimize for image denoising is the Total Variation (TV)\cite{rudin_osher_fatemi}
regularized energy function:

\[
E(u) = \lambda \| f - u \|^2_{L^2} + \int_{\Omega} \| \nabla u(x) \| \, dx
\]

where:
\begin{itemize}
    \item \( \| f - u \|^2_{L^2} = \int_{\Omega} (f(x) - u(x))^2 \, dx \) is the data fidelity term that measures the difference between the noisy image \( f \) and the clean image \( u \).
    \item \( \int_{\Omega} \| \nabla u(x) \| \, dx \) is the Total Variation (TV) norm, which enforces smoothness on the image by penalizing large gradients in \( u \).
    \item \( \lambda > 0 \) is a regularization parameter that controls the trade-off between the data fidelity and smoothness.
\end{itemize}

In a discrete setting, the objective function becomes:

\[
E(u) = \lambda \sum_{i,j} (f_{i,j} - u_{i,j})^2 + \sum_{i,j} \sqrt{(D_x u_{i,j})^2 + (D_y u_{i,j})^2}
\]

where \( D_x u_{i,j} \) and \( D_y u_{i,j} \) are the discrete gradients in the \( x \) and \( y \) directions.

\subsection{Gradient of the Objective Function}

The gradient of the objective function \( E(u) \) with respect to \( u \) is given by:

\[
\nabla_u E(u) = 2 \lambda (u - f) - \text{div} \left( \frac{\nabla u}{\| \nabla u \|} \right)
\]

where:
\begin{itemize}
    \item \( 2 \lambda (u - f) \) is the gradient of the data fidelity term.
    \item \( \text{div} \left( \frac{\nabla u}{\| \nabla u \|} \right) \) is the gradient of the Total Variation (TV) regularization term.
    \item \( \| \nabla u \| = \sqrt{(D_x u)^2 + (D_y u)^2} \) is the magnitude of the image gradient.
\end{itemize}

\subsection{Gradient Descent Algorithm}

To minimize the objective function \( E(u) \), we use the gradient descent algorithm. The update rule at iteration \( k \) is:

\[
u^{k+1} = u^k - \alpha \nabla_u E(u^k)
\]

where:
\begin{itemize}
    \item \( \alpha > 0 \) is the learning rate that controls the step size of each update.
    \item \( \nabla_u E(u^k) \) is the gradient of the objective function evaluated at \( u^k \).
\end{itemize}

At each iteration, the clean image \( u \) is updated by subtracting the gradient of the objective function scaled by the learning rate \( \alpha \).

\subsection{Results And Discussion}
Convergence Graphs, Resulting Images After Denoising and Discussion are In The Jupyter Notebook which is also attached after the report.




% References
\begin{thebibliography}{9}

\bibitem{rudin_osher_fatemi}
L. I. Rudin, S. Osher, E. Fatemi, 
\textit{Nonlinear total variation based noise removal algorithms}, 
Physica D: Nonlinear Phenomena, 
Volume 60, Issue 1-4, 1992, Pages 259-268.

\bibitem{chatgpt}
OpenAI, ChatGPT, 
\textit{ChatGPT: Language Model for Conversations}, 
OpenAI, 2023. Available at: \texttt{https://openai.com/chatgpt}.

\end{thebibliography}



\end{document}